\section{Discussion}
\label{sec:discussion}
\paragraph{Implication for sparse model design.}
The matched branches favor evaluating gate location and optimization together.
Pressure on $h$ can improve the local quality--sparsity trade-off, while
spreading the same type of pressure to four sites can increase its quality
cost. Adding attention-operand gates opens a different source of opportunity
and changes the response to pressure at the expanded site set.
This motivates architectures and objectives that specify both
\emph{which representations should be sparse} and \emph{which operations
will consume them}.

\paragraph{From elementwise zeros to executable structure.}
Our measurements suggest two directions for future interventions. Joint
activation and weight pressure could target the union of their zero operands,
with overlap explicitly accounted for. Grouped or blockwise activation
pressure could encourage channels shared across tokens, improving the
regularity available to an implementation. Such constraints introduce a
new representation trade-off: shared paths may be easier to execute but can
restrict token-specific computation. Testing them requires matched quality,
structure, and runtime measurements. The current negative execution result
provides a concrete motivation for these extensions.

\paragraph{Scope of the evidence.}
One seed, one corpus, and a fixed token budget support descriptive matched
comparisons. The threshold sweep characterizes sensitivity within this
setup; independent seeds would establish how reproducible the effects are.
The most direct next controls are A7 gates with four-site pressure and
matched Q-Sparse-style top-$k$ and optimized clipping baselines.
Downstream tasks, continued training, supervised adaptation, and
reinforcement-learning post-training would test whether the observed
trade-offs persist beyond this pretraining setting.

The 410M cohort is reported in full in Appendix~\ref{app:410m}, including
its different within-family loss response. At a common 1.493-billion-token
budget, exposure declines from 106.1 to 21.2 to 3.7 tokens per parameter
across 14M, 70M, and 410M. Compute-optimal training research motivates
increasing data together with model size \citep{hoffmann2022training};
larger, adequately budgeted comparisons therefore require more resources.
The current results leave the cause of the 410M behavior open. Data-budget
and optimization controls are needed to distinguish limited training
exposure from other explanations.
